% !PRINTED 431
\weberpage{431}
%<<CONT
sets. Equation (6) has the same \gesperrt{Galois} group as (7). We can now put any one of the roots of (7) for $U_\infty$, and consequently we can let $U_\infty$ and $v_\infty$ in (9) correspond to one another. Then we can take any one of the remaining roots of (7) for $U_0$, because we can replace $x_1, x_2$ in \S.\,113, (4) by $\varepsilon^r x_1$, $\varepsilon^{-r} x_2$. We therefore let $U_0$ and $v_0$ correspond to one another. Finally, we can still dispose of the imaginary fifth root of unity $\varepsilon$ in such a way that $U_1$ and $v_1$ correspond, and then it follows from the agreement of the group that $U_\xi$ and $v_\xi$ correspond throughout, and hence that
\weq{10}{U_\xi = -\gamma_2(\omega)v_\xi, \qquad \xi = \infty, 0, 1, 2, 3, 4}
\gesperrt{If, then, we determine} $\omega$ \gesperrt{for a given} $z$ \gesperrt{from the transcendental equation (8), then (10) and (5) give us the complete solution of the icosahedral resolvent of the sixth degree.}

By (2) and (5) we obtain, for $\xi = 0, 1, 2, 3, 4$,
\weq{11}{\eta(\omega)\sqrt{v_\xi} = \sum_{-\infty,\infty}^{\nu}(-1)^\nu e^{\frac{\pi i}{60}(\omega - 24\xi)(6\nu+1)^2}.}
The exponent $(6\nu+1)^2$ can give only the residue $0, 1, -1$ with respect to the modulus 5. If, then, we let $\nu_1, \nu_2$ run through all the numbers which satisfy the conditions
\[
\nu_1 \equiv 1, 2, \qquad \nu_2 \equiv 0, 3 \pmod{5}
\]
then we can decompose the expression (11) as follows:
\weq{12}{\eta(\omega)\sqrt{v_\xi} = \eta(5\omega) + \varepsilon^\xi \sum^{\nu_1}(-1)^{\nu_1} e^{\frac{\pi i}{60}\omega(6\nu_1+1)^2} + \varepsilon^{-\xi}\sum^{\nu_2}(-1)^{\nu_2} e^{\frac{\pi i}{60}\omega(6\nu_2+1)^2},}
if $\varepsilon = e^{\frac{2\pi i}{5}}$ is set. From this one sees that the quantities $\sqrt{v_\xi}$ satisfy the relations (9) of the preceding section, and the formula (10), \S.\,113 gives for the root $x = x_1 : x_2$ of the icosahedral equation the expression:
\weq{13}{\eta(5\omega)x = \sum^{\nu_1}(-1)^{\nu_1} e^{\frac{\pi i}{60}\omega(6\nu_1+1)^2}.}
The remaining roots are obtained by applying the icosahedral substitutions to $x$. One can, however, also
%CONT>>
